\section{Experiments and analysis}
\label{sec:experiments}
% Target: 2.5 pages. Reuse IJCAI26 evidence only; KIM results are not new results.
% Source mapping: content_mapping.md, Sections 5.1--5.5.

\subsection{Tasks, baselines, and evaluation protocol}
\label{sec:setup}
\begin{itemize}
    \item Describe low-dimensional Car Racing and Door Opening, one
    demonstration per task, demonstration provenance, and the ten reported
    unseen evaluation settings. Verify reward definitions, horizons, and
    separation between refinement rollouts and final evaluation settings.
    \item Introduce zero-shot structured policy as the primary refinement
    comparison; MLP and MLES provide context. Specify training stages and the
    different supervision/resources available to each baseline.
    \item State the LLM configuration, refinement/training budgets, metrics,
    and what each uncertainty estimate aggregates. Recover independent-run
    counts; do not equate evaluation settings with independent training runs.
\end{itemize}

\subsection{Does iterative refinement improve the initial policy?}
\label{sec:refinement-results}
\begin{itemize}
    \item Lead with the existing zero-shot versus refined Car Racing imitation
    results and the reported performance across iterations. Distinguish changes
    in fitted behavior from the untested claim of recovering the expert's strategy.
    \item Include Door Opening's reach-and-pull progression and distance metric
    as the second task. Report limited progress in intermediate iterations;
    avoid claiming that every revision produces an improvement.
    \item Plan a compact performance table and refinement-history figure using
    existing data/figures where recoverable. Keep resource columns separate
    from reward, and preserve the uncertainty definitions.
\end{itemize}

\subsection{How does the feedback affect refinement?}
\label{sec:feedback-results}
\begin{itemize}
    \item Compare generated-analysis results, subsampled trajectory tables,
    and no-feedback revision using the existing Car Racing comparison. Explain
    the table-subsampling and matched-budget details that must be recovered.
    \item Reuse the input/output token table and discuss both counts, including
    the larger output budget for analysis generation. Do not turn token counts
    into unmeasured latency, monetary cost, or total-compute savings.
    \item State the ablation's scope: it evaluates feedback procedures, not an
    isolated effect of semantic names or latent logging. No reward-only feedback
    result is present in the source manuscript.
\end{itemize}

\subsection{What changes across refinement iterations?}
\label{sec:edit-analysis}
\begin{itemize}
    \item Extend the matched method example into a compact iteration record:
    observed issue, diagnostic evidence, edit type, and subsequent behavior.
    Separate documented explanations from conjectured mechanisms. Use the
    method's running example to establish the workflow, and this subsection
    for evidence across revisions rather than retelling the same example.
    \item Reuse the existing Door Opening narrative, diagnostic categories,
    and manually coded strategy-coverage table selectively. Explain the manual
    coding procedure; move full coverage tables to the appendix.
    \item Include ineffective revisions and diagnostic blind spots where
    documented. Recover per-iteration programs/logs to distinguish additions,
    operation/constraint changes, and initialization changes.
\end{itemize}

\subsection{Does refinement help subsequent reinforcement learning?}
\label{sec:rl-results}
\begin{itemize}
    \item Present the reported refined versus zero-shot policy outcomes at the
    stated RL budgets as a downstream benefit. Briefly explain when RL begins;
    keep its detailed derivation in Appendix~\ref{app:rl}.
    \item Separate refinement resets, RL resets, and LLM calls. Audit the old
    abstract's compute percentage against the table; report environment resets
    as such, without claiming measured total-compute savings or statistical
    equivalence from similar mean returns.
\end{itemize}
